\section{The sharp four-state amplitude bound under all triplet tests}
\label{sec:fulltripletbound}
The additional $F$ inequalities sharpen the contrast-cone bound from one
half to one quarter. The key restriction concerns the first three entries
of the slow eigenfunction, rather than monotonicity of all four entries.

\begin{lemma}[First-three arithmetic restriction]\label{lem:firstthreearithmetic}
Let a positive reversible four-state generator have distinct ordered
energies, a finite positive bath, a simple slow mode, and $F_{ijk}\leq0$,
$G_{ijk}\leq0$ on every original-rate triplet. Orient its slow left
eigenfunction by the nonnegative contrast coordinates of
Theorem~\ref{thm:orderedk}. Then
\[
w_1\geq w_2,\qquad w_1+w_2\geq2w_3,\qquad w_4\leq0.
\]
\end{lemma}
\begin{proof}
Write $q_i=\pi_i(\beta_b)$, $Q_j=\sum_{i\leq j}q_i$, and use the
unnormalized prefix contrasts $h^{(1)},h^{(2)},h^{(3)}$ from the contrast
theorem. Their coordinate generator is
\[
B=\operatorname{diag}(\eta_1,\eta_2,\eta_3)^{-1}H^\top D L H,
\quad H=(h^{(1)},h^{(2)},h^{(3)}),\quad D=\operatorname{diag}(q).
\]
It is diagonally similar to the symmetric Metzler contrast matrix. Its
Perron eigenvector $x$ can be chosen nonnegative and nonzero, and $w=Hx$.
Consequently $w_1-w_2=(Q_2/q_2)x_1\geq0$ and
$w_4=-(Q_3/q_4)x_3\leq0$. Put
\[
s=\frac{(q_1-q_2)q_3}{2q_2Q_3}>0.
\]
The remaining assertion is equivalent to $x_2\geq sx_1$, since
$w_1+w_2-2w_3=(2Q_3/q_3)(x_2-sx_1)$.

We first make a reset subtraction, which is used only in this proof and
does not alter the original triplet tests. The reduced conductances are
$k_{ij}=a_{ij}/q_i$. Their minimum is in the first row, by all the $G$
inequalities. Subtract
$k_* (q\mathbf1^\top-Q_4 I)$, where
$k_* =\min\{k_{12},k_{13},k_{14}\}$.
All resulting rates are nonnegative; their generator is reversible and
conservative, although it may be reducible. Every $F$ and $G$ is unchanged.
On the nonstationary contrast space this operation adds $k_*Q_4I$ to $B$,
so its eigenvectors, simplicity and eigenvalue ordering are unchanged.
The shifted largest contrast eigenvalue, denoted $\lambda$, is at most zero,
because the shifted full matrix remains a Markov generator. At least one
of the shifted downward rates $a_{12},a_{13},a_{14}$ is zero. All quantities
in the rest of this proof refer to this shifted generator.

If $x_1=0$, the assertion is immediate. Suppose $x_1>0$. Then
$\delta=\lambda-B_{33}>0$. To see strictness, symmetrize $B$: its largest
eigenvalue is at least its third diagonal entry. Equality would force the
third row and column of the positive semidefinite matrix
$\lambda I-\operatorname{diag}(\sqrt\eta)B
\operatorname{diag}(1/\sqrt\eta)$ to vanish. The third coordinate vector
would then be a Perron eigenvector. Simplicity would force $x_1=0$, a
contradiction. This also handles a zero shifted eigenvalue.

Eliminate $x_3$ from the eigenvector equations. The resulting two by two
matrix is
\[
\mathcal B_{ij}=B_{ij}+\frac{B_{i3}B_{3j}}{\delta},\qquad i,j\leq2.
\]
It has nonnegative off-diagonal entries and is diagonally symmetrizable.
The Schur complement of $\lambda I-B$ shows that $\lambda$ is its simple
largest eigenvalue, with eigenvector $(x_1,x_2)$.
Define
\[
C=B_{21}+s(B_{22}-B_{11})-s^2B_{12},\qquad
D_*=(B_{23}-sB_{13})(B_{31}+sB_{32}).
\]
Then
\[
\delta\{\mathcal B_{21}+s(\mathcal B_{22}-\mathcal B_{11})
-s^2\mathcal B_{12}\}=\delta C+D_*.
\]
A nonnegative value of this expression implies $x_2/x_1\geq s$. Indeed,
if the two off-diagonal entries of $\mathcal B$ are positive, its Perron
ratio is the unique positive root of the concave quadratic in braces,
viewed as a function of $s$. If those entries vanish, symmetrizability makes
both vanish. Since $x_1>0$ and the largest eigenvalue is simple, one then
has $x_2=0$ and $\mathcal B_{11}>\mathcal B_{22}$; the expression in braces
is strictly negative. Thus this exceptional case cannot occur when the
expression is nonnegative.

All $G$ inequalities give $B_{31},B_{32}\geq0$, and direct simplification gives
\[
B_{23}-sB_{13}
=\frac{Q_4}{2q_1q_2Q_3}
\{2a_{34}q_1q_2-a_{14}q_2q_3-a_{24}q_1q_3\}\geq0.
\]
Hence $D_*\geq0$. If $C\geq0$ the desired expression is already
nonnegative. If $C<0$, use $\lambda\leq0$ to obtain
\[
\delta C+D_*\geq-B_{33}C+D_*.
\]
It remains to prove that this last, eigenvalue-free expression is nonnegative.

For brevity write the shifted downward rates as
$(a,b,c,d,e,f)=(a_{12},a_{13},a_{14},a_{23},a_{24},a_{34})$ and define
\[
N=-2q_1q_2q_3F_{123}\geq0,\qquad
J=e q_1(2q_2+q_3)-c q_2(2q_1+q_3).
\]
Exact multiplication yields
\[
-B_{33}C+D_*=
\frac{Q_2Q_4}{4q_1q_2^2Q_3^2}\,T,
\quad
T=(c+e+f)N+q_4\{2ceq_3(q_1-q_2)+fJ\}.
\]
The prefactor is positive. Since at least one of $a,b,c$ is zero,
nonnegativity follows from these three exhaustive cases:
\begin{enumerate}
\item If $c=0$, then $J=e q_1(2q_2+q_3)\geq0$ and all terms of $T$
are nonnegative.
\item If $a=0$, $F_{124}\leq0$ gives
$c q_2(2q_1+q_4)\leq e q_1(2q_2+q_4)$.
The ratio $(2q_2+z)/(2q_1+z)$ strictly increases with $z$ and
$q_3>q_4$. It follows, without division by $c$ or $e$, that $J\geq0$.
Thus $T\geq0$.
\item If $b=0$, then
$N=2a q_2q_3(q_1-q_2)+d q_1q_3(2q_2+q_3)$.
If $c\leq e$, one has $J\geq0$. If $c>e$, the inequality
$F_{124}\leq0$ gives instead
\[
N+q_4J\geq
2q_1q_2(q_3-q_4)(c-e)+d q_1q_3(2q_2+q_3)\geq0.
\]
Now write
$T=(c+e)N+f(N+q_4J)+2ceq_3q_4(q_1-q_2)\geq0$.
\end{enumerate}
This proves $x_2\geq sx_1$ and completes the lemma, including equality
faces and every reducible reset-subtracted case.
\end{proof}

\begin{theorem}[Sharp quarter-amplitude bound]\label{thm:fulltripletbound}
Under the hypotheses of Lemma~\ref{lem:firstthreearithmetic}, for every
$0<\beta_h<\beta_c<\beta_b$,
\[
\frac{|b_2(\beta_c)|}{|b_2(\beta_h)|}<\frac43.
\]
Consequently every ordinary cooling amplitude advantage is strictly below
$25\%$. The bound is sharp as a supremum, by
Theorem~\ref{thm:quarterlimit}. The corresponding supremum of the
asymptotic relative-entropy advantage is $7/16$, and the time advance is
strictly below $\log(4/3)/|\lambda_2|$.
\end{theorem}
\begin{proof}
Let $m(\beta)$ be the conditional Gibbs expectation of $w$ on states
$1,2,3$. It is nondecreasing for $\beta\geq0$. To verify this directly,
shift $E_1=0$, put $a=E_2$, $b=E_3$, $t=e^{-\beta a}$ and
$y=e^{-\beta b}$, and use conditional probabilities
$(p_1,p_2,p_3)=(1,t,y)/(1+t+y)$. The first-three vector is a constant
plus
\[
u(1,-1,0)+v(1,1,-2),\qquad
u=(w_1-w_2)/2\geq0,\quad v=(w_1+w_2-2w_3)/6\geq0.
\]
Furthermore
\[
(p_1-p_2)'=\frac{2at+by-(b-a)ty}{(1+t+y)^2}>0,
\qquad (p_1+p_2-2p_3)'=-3p_3'>0.
\]
The first inequality uses $t\leq1$ and $b>a>0$; the second uses that
$E_3$ is the largest of the three energies. Thus $m$ is nondecreasing,
strictly unless the first three mode components coincide.

Let $P(\beta)=Z_3(\beta)/Z_4(\beta)$ be the Gibbs mass on the first three
states. It strictly increases with $\beta$, since the excluded state has
strictly higher energy. For finite positive $\beta$, $3/4<P(\beta)<1$.
Write $P_h=P(\beta_h)$, $P_c=P(\beta_c)$, and $r=P_c/P_h<4/3$.
The full overlap is $b_2=Pm+(1-P)w_4$. Therefore
\[
b_2(\beta_c)-r b_2(\beta_h)
=P_c\{m(\beta_c)-m(\beta_h)\}+(1-r)w_4\geq0.
\]
The contrast theorem ensures both overlaps are strictly negative on cooling,
so $|b_2(\beta_c)|\leq r|b_2(\beta_h)|<(4/3)|b_2(\beta_h)|$.
The relative-entropy and time-advance statements follow from their respective
quadratic and logarithmic dependence on the amplitude ratio. The existing
quarter-limit family proves sharpness; none of these suprema is attained at
finite parameters in the stated domain.
\end{proof}
